\section*{Chapter \ref{ch:g2:from-seed-to-plant} --- From Seed to Plant}

\begin{solution}{exo:g2:from-seed-to-plant:1}
A tiny \omterm{def:g1:plants-around-us:plant}{plant} already started (the germ of the \omterm{def:g1:plants-around-us:plant}{plant}, with the
beginnings of a \omterm{def:g1:plants-around-us:parts}{root} and a shoot) and a store of food to feed it.
\end{solution}

\begin{solution}{exo:g2:from-seed-to-plant:2}
To \omterm{def:g2:from-seed-to-plant:germination}{germinate} is to wake and start growing: the \omterm{def:g2:from-seed-to-plant:seed}{seed} swells, its
coat splits, and it begins to become a \omterm{def:g1:plants-around-us:plant}{plant}. The \omterm{def:g1:plants-around-us:parts}{root} comes out
first, and dives downward.
\end{solution}

\begin{solution}{exo:g2:from-seed-to-plant:3}
Water and warmth. Light is not on the list --- the \omterm{def:g2:from-seed-to-plant:seed}{seed} \omterm{def:g2:from-seed-to-plant:germination}{germinates}
happily in the dark, under the soil.
\end{solution}

\begin{solution}{exo:g2:from-seed-to-plant:4}
It lives off the food store packed inside the \omterm{def:g2:from-seed-to-plant:seed}{seed} --- the two
thick halves of the bean. That store feeds the \omterm{prop:g2:from-seed-to-plant:cycle}{seedling} until its
green \omterm{def:g1:plants-around-us:parts}{leaves} reach the light and it can feed itself.
\end{solution}

\begin{solution}{exo:g2:from-seed-to-plant:5}
For example: beans, peas, lentils, rice, sunflower \omterm{def:g2:from-seed-to-plant:seed}{seeds}, apple
pips, a peach stone.
\end{solution}

\begin{solution}{exo:g2:from-seed-to-plant:6}
Because a \omterm{def:g2:from-seed-to-plant:seed}{seed} needs warmth as well as water: in cold winter earth
it would wait, or rot, instead of germinating. Spring brings the
warm soil the \omterm{def:g2:from-seed-to-plant:seed}{seed} is waiting for.
\end{solution}

\begin{solution}{exo:g2:from-seed-to-plant:7}
\omterm{def:g2:from-seed-to-plant:seed}{Seed}, \omterm{prop:g2:from-seed-to-plant:cycle}{seedling}, leafy \omterm{def:g1:plants-around-us:plant}{plant}, flowering \omterm{def:g1:plants-around-us:plant}{plant}, pod with new beans.
\end{solution}

\begin{solution}{exo:g2:from-seed-to-plant:8}
Both \omterm{def:g1:plants-around-us:parts}{roots} grow downward and both shoots grow upward, whichever
way the beans were lying. The \omterm{prop:g2:from-seed-to-plant:cycle}{seedling} always finds up and down
--- exactly what the remark predicts.
\end{solution}

\begin{solution}{exo:g2:from-seed-to-plant:9}
In the drawer the \omterm{def:g2:from-seed-to-plant:seed}{seeds} had no water, so they could not \omterm{def:g2:from-seed-to-plant:germination}{germinate};
but a good \omterm{def:g2:from-seed-to-plant:seed}{seed} without water does not die --- it waits. A year
later, given water and warmth, the waiting \omterm{def:g2:from-seed-to-plant:seed}{seeds} woke normally.
\end{solution}

\begin{solution}{exo:g2:from-seed-to-plant:10}
The \omterm{def:g2:from-seed-to-plant:seed}{seed}'s food store --- the bean's two thick halves --- plays
the yolk's part. \omterm{def:g2:animals-grow-up:born}{Egg} and \omterm{def:g2:from-seed-to-plant:seed}{seed} are both beginnings of a new \omterm{def:g1:living-or-not:living}{living
thing}, packed with the food the young one needs before it can
feed itself.
\end{solution}
